\section{Conclusion}

A self-replicating space factory cannot close on its own chips, and the reason is not a
gap that better engineering will close. Chips sit at the end of the deepest supply chain on
Earth and cost about three orders of magnitude more energy per kilogram to make than the
structure that houses them, so a lone seed that tries to make them either fails to
reproduce the pipeline or runs out of power first. That cap on closure is what bounds the
launch-mass leverage $1/(1-C)$ that motivates the whole concept: the last points of closure
buy most of the payoff, and the electronics wall holds a factory short of them. Local
chip-making rewards only a seed that is both highly closed and richly powered, and backfires
otherwise. The design target that follows is a factory closed on nearly everything except
electronics, fed by imported vitamins, and sized so its power budget never tempts it into
making the one part it cannot afford. The leverage of self-replication is real; the
electronics wall is what bounds it, and knowing where the wall stands is what lets a seed be
designed to live with it rather than break against it.

\section*{Acknowledgment}
This work used no external funding. The figures and reported figures-of-merit regenerate
deterministically from the project's closure model, and every quantitative claim traces to
a cited source in the project's shared bibliography.
